\begin{lemma}
Under the hypotheses of \(\noderef{HugeFamilySizeBound}\), and assuming also
\(\varepsilon_1\ge0\), suppose the sample parameters are the rounded hierarchy
parameters
\[
m=\noderef{TwoBiteNaturalReducedVertexCount}(n),
\qquad
p=\noderef{TwoBiteEdgeProbability}\!\left(\frac12,n\right).
\]
Then the same-color huge projected-pair sums satisfy
\[
\sum_{x\in H_I\cap V_R}\binom{|\pi_R(X_x(I))|}{2}_{\mathbb R}
\le\varepsilon_1k^2,
\]
\[
\sum_{x\in H_I\cap V_B}\binom{|\pi_B(X_x(I))|}{2}_{\mathbb R}
\le\varepsilon_1k^2,
\]
and their sum is at most \(\varepsilon_1k^2\).
\end{lemma}
\begin{proof}
Let
\[
H_I=\{x:\noderef{TwoBiteHugeClass}(\omega,I,x)\}
\]
be the full huge filter.  In the Lean statement the two same-color filters are
the finite filtered-universe sets attached to the predicates
\[
\mathsf{hugeRed}(x)\iff
\noderef{TwoBiteHugeClass}(\omega,I,x)\wedge
  \exists r,\ x=\operatorname{inl}(r),
\]
and
\[
\mathsf{hugeBlue}(x)\iff
\noderef{TwoBiteHugeClass}(\omega,I,x)\wedge
  \exists b,\ x=\operatorname{inr}(b).
\]
Thus the exact finite filters are
\[
H_R=\operatorname{univ}.\operatorname{filter}(\mathsf{hugeRed}),
\qquad
H_B=\operatorname{univ}.\operatorname{filter}(\mathsf{hugeBlue}).
\]
Equivalently, \(H_R=H_I\cap V_R\) and \(H_B=H_I\cap V_B\).  Membership in
\(H_R\) gives both hugeness and a representation \(x=\operatorname{inl}(r)\)
for a unique red base vertex \(r\); membership in \(H_B\) gives
\(x=\operatorname{inr}(b)\) for a unique blue base vertex \(b\).
The hypotheses \(m=\noderef{TwoBiteNaturalReducedVertexCount}(n)\) and
\(p=\noderef{TwoBiteEdgeProbability}(\frac12,n)\) identify the sample
parameters with the rounded parameters in
\(\noderef{ParameterHierarchyEventualComparisons}\).  Thus all hierarchy
comparisons below involving \(2pm\), including \(1\le 2pm\) and the
\((2k/t_1)\binom{2pm}{2}_{\mathbb R}\) bound, are being applied to the same
\(m\) and \(p\) that occur in the sample \(\omega\).
Unfold \(\noderef{TwoBiteProjectedPairSum}\).  The red same-color quantity is
\[
S_R=\sum_{x\in H_R}
  \binom{
    |(\noderef{TwoBiteX}(\omega,I,x)).\mathrm{image}
      (\noderef{TwoBiteRedProjection}(\omega))|
  }{2}_{\mathbb R},
\]
and the blue same-color quantity is
\[
S_B=\sum_{x\in H_B}
  \binom{
    |(\noderef{TwoBiteX}(\omega,I,x)).\mathrm{image}
      (\noderef{TwoBiteBlueProjection}(\omega))|
  }{2}_{\mathbb R}.
\]
Here and below the shorter notation \(|\pi_R(X_x(I))|\) and
\(|\pi_B(X_x(I))|\) denotes these exact real cardinalities of finite images.
Using the exact predicates and projections from the Lean statement,
\[
\noderef{TwoBiteProjectedPairSum}(\omega,I,\mathsf{hugeRed},\pi_R)=S_R,
\qquad
\noderef{TwoBiteProjectedPairSum}(\omega,I,\mathsf{hugeBlue},\pi_B)=S_B.
\]
Unfolding \(\noderef{ClosedPairsBound}\), the three conjuncts in the Lean
statement are exactly
\[
S_R\le\varepsilon_1k^2,\qquad
S_B\le\varepsilon_1k^2,\qquad
S_R+S_B\le\varepsilon_1k^2.
\]

We first record the pointwise projection-size estimate.  Fix \(x\in H_R\).
Then \(x=\operatorname{inl}(r)\) for a unique red base vertex \(r\).  If
\(z\in X_x(I)\), then by unfolding \(\noderef{TwoBiteX}\),
\[
z\in I\quad\hbox{and}\quad z\in N_{\operatorname{inl}(r)}.
\]
The lifted-neighborhood definition used by \(\noderef{TwoBiteX}\) says that
\[
z\in N_{\operatorname{inl}(r)}
\quad\Longleftrightarrow\quad
G_R.Adj(r,\noderef{TwoBiteRedProjection}(\omega,z)).
\]
Let
\[
N_R(r)=\{u:\ G_R.Adj(r,u)\}
\]
be the finite red base-neighbor set; its cardinality is the red degree count
appearing in the base-degree clause of \(\noderef{FiberAndDegreeEvent}\).
Concretely, this is the filtered-universe set
\[
\operatorname{univ}.\operatorname{filter}
  \bigl(u\mapsto G_R.Adj(r,u)\bigr),
\]
so its cardinality is
\[
\operatorname{GraphDegreeCount}(G_R,r).
\]
Thus the red base-degree conjunct of \(\noderef{FiberAndDegreeEvent}\) is
being applied to exactly the real cardinality \((|N_R(r)|:\mathbb R)\), with
target \(pm\).
To see the containment in the finite-image form used by Lean, take
\[
u\in(\noderef{TwoBiteX}(\omega,I,x)).\operatorname{image}
      (\noderef{TwoBiteRedProjection}(\omega)).
\]
Then \(u=\noderef{TwoBiteRedProjection}(\omega,z)\) for some \(z\in X_x(I)\).
The preceding equivalence applied to this \(z\), followed by substitution of
the displayed equality for \(u\), gives \(u\in N_R(r)\).  Hence
\[
(\noderef{TwoBiteX}(\omega,I,x)).\operatorname{image}
      (\noderef{TwoBiteRedProjection}(\omega))
\subseteq N_R(r). \tag{1}
\]
Taking finite cardinalities in (1), and then casting to \(\mathbb R\), the
image-cardinality is at most \(|N_R(r)|\).  The red base-degree clause of
\(\noderef{FiberAndDegreeEvent}\), after unpacking the relative-error
condition, gives
\[
\bigl||N_R(r)|-pm\bigr|\le\varepsilon_2pm.
\]
The upper side of this absolute-value bound gives
\[
|N_R(r)|-pm\le\varepsilon_2pm,
\]
so
\[
|N_R(r)|\le pm+\varepsilon_2pm=(1+\varepsilon_2)pm.
\]
The conjunct \(0\le\varepsilon_2\le1\) of
\(\noderef{ParameterHierarchyEventualComparisons}\), together with the rounded
comparison \(1\le2pm\), gives \(0\le pm\).  Therefore
\(\varepsilon_2pm\le pm\), and
\[
|(\noderef{TwoBiteX}(\omega,I,x)).\operatorname{image}
      (\noderef{TwoBiteRedProjection}(\omega))|
\le2pm. \tag{2}
\]

The blue estimate is the same argument with colors exchanged.  If \(x\in H_B\)
then \(x=\operatorname{inr}(b)\) for a unique blue base vertex \(b\).  For
each \(z\in X_x(I)\), unfolding \(\noderef{TwoBiteX}\) and the blue branch of
the lifted-neighborhood definition gives
\[
G_B.Adj(b,\noderef{TwoBiteBlueProjection}(\omega,z)).
\]
Let \(N_B(b)=\{u:\ G_B.Adj(b,u)\}\).  Its cardinality is the blue degree count
appearing in \(\noderef{FiberAndDegreeEvent}\): in Lean it is the cardinality
of
\[
\operatorname{univ}.\operatorname{filter}
  \bigl(u\mapsto G_B.Adj(b,u)\bigr),
\]
namely \(\operatorname{GraphDegreeCount}(G_B,b)\).  Thus the blue
base-degree conjunct applies to the exact real cardinality
\((|N_B(b)|:\mathbb R)\), again with target \(pm\).
Thus an arbitrary element \(u\) of the finite image
\((\noderef{TwoBiteX}(\omega,I,x)).\operatorname{image}
      (\noderef{TwoBiteBlueProjection}(\omega))\)
has the form \(u=\noderef{TwoBiteBlueProjection}(\omega,z)\) with
\(z\in X_x(I)\), and is a blue base-neighbor of \(b\).  Hence
\[
(\noderef{TwoBiteX}(\omega,I,x)).\operatorname{image}
      (\noderef{TwoBiteBlueProjection}(\omega))
\subseteq N_B(b).
\]
Taking finite cardinalities and casting to \(\mathbb R\), the image size is at
most \(|N_B(b)|\).  The blue base-degree clause of
\(\noderef{FiberAndDegreeEvent}\), the same relative-error calculation,
\(0\le\varepsilon_2\le1\), and \(1\le2pm\) give
\[
\bigl||N_B(b)|-pm\bigr|\le\varepsilon_2pm,\qquad
|N_B(b)|\le(1+\varepsilon_2)pm\le2pm.
\]
Together with the image containment, this gives
\[
|(\noderef{TwoBiteX}(\omega,I,x)).\operatorname{image}
      (\noderef{TwoBiteBlueProjection}(\omega))|
\le2pm. \tag{3}
\]

Let \(B=2pm\).  The lower-bound conjunct of
\(\noderef{ParameterHierarchyEventualComparisons}\), evaluated at the present
\(n>n_0\) with the identified rounded parameters
\[
m=\noderef{TwoBiteNaturalReducedVertexCount}(n),\qquad
p=\noderef{TwoBiteEdgeProbability}\!\left(\frac12,n\right),
\]
gives
\[
1\le2pm=B. \tag{B}
\]
For each red or blue
summand, let
\[
a=|\pi_R(X_x(I))|\quad\hbox{or}\quad a=|\pi_B(X_x(I))|.
\]
Then \(a\) is the real cast of a natural number \(c\), the corresponding image
cardinality, so \(a\ge0\), and (2) or (3) gives \(a\le B\).  We now justify
the binomial comparison for this actual cardinality.  By
\(\noderef{RealChooseTwo}\),
\[
\binom{u}{2}_{\mathbb R}=\frac{u(u-1)}2.
\]
For \(0\le u\le v\) with \(1\le u\), the difference
\[
\binom{v}{2}_{\mathbb R}-\binom{u}{2}_{\mathbb R}
  =\frac{(v-u)(u+v-1)}2
\]
is nonnegative; hence the function is increasing on \([1,\infty)\).  If
\(c\ge2\), then \(1\le a\), so this monotonicity and \(a\le B\) give
\(\binom a2_{\mathbb R}\le\binom B2_{\mathbb R}\).  If \(c=0\), then
\(a=0\) and \(\binom a2_{\mathbb R}=0\).  If \(c=1\), then \(a=1\) and again
\(\binom a2_{\mathbb R}=0\).  In these two exceptional cardinality cases,
\(\binom B2_{\mathbb R}\ge0\), because \(1\le B\) implies
\[
\binom B2_{\mathbb R}=\frac{B(B-1)}2\ge0.
\]
Thus in all cardinality cases
\[
\binom{a}{2}_{\mathbb R}\le \binom{B}{2}_{\mathbb R}. \tag{4}
\]
Applying (4) to all red and blue summands yields
\[
S_R+S_B
\le (|H_R|+|H_B|)\binom{2pm}{2}_{\mathbb R}. \tag{5}
\]
Indeed, for the red filter, finite-sum monotonicity applies over the exact
finite domain \(H_R\) with the pointwise comparison (4), giving
\[
S_R\le (|H_R|:\mathbb R)\binom{2pm}{2}_{\mathbb R}
\]
because the upper summand is the constant
\(\binom{2pm}{2}_{\mathbb R}\).  The finite sum of that constant over \(H_R\)
is the real cast of \(|H_R|\) times the constant.  The same argument over
\(H_B\) gives
\[
S_B\le (|H_B|:\mathbb R)\binom{2pm}{2}_{\mathbb R},
\]
and adding these two inequalities gives (5).
The filters \(H_R\) and \(H_B\) are disjoint: if a vertex lay in both, then
there would be \(r,b\) with
\(\operatorname{inl}(r)=\operatorname{inr}(b)\), impossible for the two
constructors of the disjoint sum.  Each filter is also a subfamily of the full
huge filter \(H_I=\{x:\noderef{TwoBiteHugeClass}(\omega,I,x)\}\).  Therefore
the finite-cardinality identity for disjoint unions gives
\[
|H_R\cup H_B|=|H_R|+|H_B|,
\]
and monotonicity of cardinality under the inclusion
\(H_R\cup H_B\subseteq H_I\) gives
\[
|H_R|+|H_B|=|H_R\cup H_B|\le |H_I|. \tag{6}
\]
The quantity \(\binom{2pm}{2}_{\mathbb R}\) is nonnegative by (B).
Multiplying (6), after casting finite cardinalities to real numbers, by this
nonnegative factor and using (5), we obtain
\[
S_R+S_B
\le |H_I|\binom{2pm}{2}_{\mathbb R}. \tag{7}
\]

Apply the second conclusion of \(\noderef{HugeFamilySizeBound}\) with the
same sample \(\omega\), set \(I\), integer \(k\), integer \(n\), parameters
\(\varepsilon,\varepsilon_1,\varepsilon_2,n_0\), the same
\(\noderef{ParameterHierarchyEventualComparisons}\) hypothesis, the same
strict inequality \(n_0<n\), the same equality
\[
k=\noderef{TwoBiteNaturalIndependenceScale}(1+\varepsilon,n),
\]
the same cardinality hypothesis \(|I|=k\), and the same event
\(\noderef{FiberAndDegreeEvent}(\omega,\varepsilon_2)\).  These are exactly
the hypotheses of \(\noderef{HugeFamilySizeBound}\).  Its full huge filter is
precisely \(H_I\), and with \(t_1=\noderef{TwoBiteHugeCutoff}(n)\) it gives
\[
|H_I|\le\frac{2k}{t_1}. \tag{8}
\]
The factor \(\binom{2pm}{2}_{\mathbb R}\) is nonnegative by (B) and the
definition of \(\noderef{RealChooseTwo}\).  Apply the elementary numerical
step \(\noderef{HugeSameColorProjectionPairNumerics}\) with
\[
H=|H_I|,\qquad t_1=\noderef{TwoBiteHugeCutoff}(n).
\]
Its three inputs are exactly (8), the lower bound (B), and the same-color
huge-projection pair-count comparison from
\(\noderef{ParameterHierarchyEventualComparisons}\), evaluated at the present
\(n>n_0\).  It gives
\[
|H_I|\binom{2pm}{2}_{\mathbb R}
\le
\frac{2k}{t_1}\binom{2pm}{2}_{\mathbb R}
\le\varepsilon_1k^2. \tag{9}
\]
The last displayed inequality is exactly the conjunct
\[
\frac{2k}{t_1}\binom{2pm}{2}_{\mathbb R}\le\varepsilon_1k^2
\]
from \(\noderef{ParameterHierarchyEventualComparisons}\), with the rounded
scales \(k=K\), \(m=M\), \(p=\frac12\sqrt{\log n/n}\), and
\(t_1=\noderef{TwoBiteHugeCutoff}(n)\).  Combining (7) and (9) proves
\[
S_R+S_B\le\varepsilon_1k^2. \tag{10}
\]

It remains only to separate the two individual bounds.  Each summand in
\(S_R\) and \(S_B\) is nonnegative.  To check this in the form used by Lean,
let the finite image cardinality be the natural number \(c\), and put
\(a=(c:\mathbb R)\).  If \(c=0\), then \(a=0\) and
\(\binom a2_{\mathbb R}=0\).  If \(c=1\), then \(a=1\) and again
\(\binom a2_{\mathbb R}=0\).  If \(c\ge2\), then \(a\ge1\), so both \(a\)
and \(a-1\) are nonnegative, and
\[
\binom a2_{\mathbb R}=\frac{a(a-1)}2\ge0.
\]
This exhaustive natural-cardinality case split proves every summand in the
two finite sums is nonnegative.  Finite sums of nonnegative real summands are
nonnegative, so \(0\le S_R\) and \(0\le S_B\).  Using these inequalities and
(10),
\[
S_R\le S_R+S_B\le\varepsilon_1k^2
\]
and
\[
S_B\le S_R+S_B\le\varepsilon_1k^2.
\]
Together with the already proved combined inequality, and after unfolding
\(\noderef{ClosedPairsBound}\), these are exactly the red, blue, and combined
closed-pair bounds asserted in the statement.
\end{proof}
